\documentclass[12pt, a4paper]{ctexart}
\usepackage[margin=2cm]{geometry}
\usepackage{amsmath, amssymb}
\usepackage{enumitem}
\usepackage{fancyhdr}
\usepackage{graphicx}
\usepackage{multicol}
\usepackage{array}
\usepackage{tabularx}
\usepackage{tikz}
\usepackage{CJKfntef}

\setCJKmainfont{Hiragino Sans GB}
\setCJKsansfont{Heiti SC}

\pagestyle{fancy}
\fancyhf{}
\fancyhead[L]{七年级上册数学测试卷}
\fancyhead[R]{第~\thepage~页}
\renewcommand{\headrulewidth}{0.4pt}
\fancyfoot[C]{共~4~页}

\setlength{\parindent}{0pt}
\setlength{\parskip}{6pt}

\newcommand{\blank}[1]{\underline{\hspace{#1}}}

\begin{document}

\begin{center}
{\LARGE\bfseries 七年级上册数学测试卷}\\[8pt]
{\large 满分：120分\quad 考试时间：120分钟}\\[4pt]
{\normalsize 命题范围：有理数、整式的加减、一元一次方程、几何初步}
\end{center}

\vspace{4pt}
\noindent\textbf{注意事项：}
\begin{enumerate}[leftmargin=2em]
\item 本试卷共三大题，24小题。
\item 选择题每小题3分，填空题每小题3分，解答题共72分。
\item 请在答题卡上作答，在本试卷上作答无效。
\end{enumerate}

\vspace{6pt}
\noindent\rule{\textwidth}{0.8pt}

\vspace{8pt}
\noindent{\large\bfseries 一、选择题（本大题共10小题，每小题3分，共30分。每小题给出的四个选项中，只有一个是正确的。）}

\vspace{6pt}

\begin{enumerate}[label=\textbf{\arabic*.}, leftmargin=2.5em]

\item $-5$ 的绝对值是（\blank{1cm}）
\begin{enumerate}[label=\Alph*., nosep]
\item $5$
\item $-5$
\item $\dfrac{1}{5}$
\item $-\dfrac{1}{5}$
\end{enumerate}

\vspace{6pt}

\item 计算 $(-3) \times 4$ 的结果是（\blank{1cm}）
\begin{enumerate}[label=\Alph*., nosep]
\item $12$
\item $-12$
\item $7$
\item $-7$
\end{enumerate}

\vspace{6pt}

\item 在数轴上，表示 $-2$ 的点在表示 $1$ 的点的（\blank{1cm}）
\begin{enumerate}[label=\Alph*., nosep]
\item 左边
\item 右边
\item 上方
\item 无法确定
\end{enumerate}

\vspace{6pt}

\item 下列运算正确的是（\blank{1cm}）
\begin{enumerate}[label=\Alph*., nosep]
\item $3a + 2b = 5ab$
\item $3a^2 - a^2 = 2$
\item $5a^2b - 3a^2b = 2a^2b$
\item $a + a = a^2$
\end{enumerate}

\vspace{6pt}

\item 化简 $2(x - 3) + 3(x + 1)$ 的结果是（\blank{1cm}）
\begin{enumerate}[label=\Alph*., nosep]
\item $5x - 3$
\item $5x + 3$
\item $5x - 9$
\item $5x + 9$
\end{enumerate}

\vspace{6pt}

\item 方程 $2x + 5 = 11$ 的解是（\blank{1cm}）
\begin{enumerate}[label=\Alph*., nosep]
\item $x = 2$
\item $x = 3$
\item $x = 4$
\item $x = 8$
\end{enumerate}

\vspace{6pt}

\item 若 $x = -1$ 是方程 $ax + 3 = 1$ 的解，则 $a$ 的值是（\blank{1cm}）
\begin{enumerate}[label=\Alph*., nosep]
\item $2$
\item $-2$
\item $4$
\item $-4$
\end{enumerate}

\vspace{6pt}

\item 下面图形中，属于棱柱的是（\blank{1cm}）
\begin{enumerate}[label=\Alph*., nosep]
\item 圆柱
\item 圆锥
\item 长方体
\item 球
\end{enumerate}

\vspace{6pt}

\item 如图，点 $A$、$B$、$C$ 在同一条直线上，若 $AB = 5$，$BC = 3$，则 $AC$ 的长为（\blank{1cm}）
\begin{enumerate}[label=\Alph*., nosep]
\item $2$
\item $8$
\item $2$ 或 $8$
\item $15$
\end{enumerate}

\vspace{6pt}

\item 某天早晨的气温是 $-3\,^{\circ}\mathrm{C}$，中午上升了 $8\,^{\circ}\mathrm{C}$，到傍晚又下降了 $4\,^{\circ}\mathrm{C}$，则傍晚的气温是（\blank{1cm}）
\begin{enumerate}[label=\Alph*., nosep]
\item $1\,^{\circ}\mathrm{C}$
\item $-1\,^{\circ}\mathrm{C}$
\item $5\,^{\circ}\mathrm{C}$
\item $7\,^{\circ}\mathrm{C}$
\end{enumerate}

\end{enumerate}

\vspace{10pt}
\noindent\rule{\textwidth}{0.8pt}

\vspace{8pt}
\noindent{\large\bfseries 二、填空题（本大题共6小题，每小题3分，共18分。）}

\vspace{6pt}

\begin{enumerate}[label=\textbf{\arabic*.}, leftmargin=2.5em, resume]

\item 比较大小：$-3$ \blank{1.5cm} $-5$（填"$>$"或"$<$"）。

\vspace{6pt}

\item 计算：$(-2)^3 = $ \blank{2cm} 。

\vspace{6pt}

\item 单项式 $-\dfrac{3x^2y}{5}$ 的系数是 \blank{2cm} ，次数是 \blank{2cm} 。

\vspace{6pt}

\item 若 $|a| = 4$，$|b| = 2$，且 $a < 0$，$b > 0$，则 $a - b = $ \blank{2cm} 。

\vspace{6pt}

\item 一个角的补角比它的余角大 $60^{\circ}$，则这个角的度数为 \blank{2cm} 。

\vspace{6pt}

\item 如图，$O$ 是直线 $AB$ 上一点，$\angle AOC = 50^{\circ}$，$OD$ 平分 $\angle BOC$，则 $\angle BOD = $ \blank{2cm} 。

\end{enumerate}

\vspace{10pt}
\noindent\rule{\textwidth}{0.8pt}

\vspace{8pt}
\noindent{\large\bfseries 三、解答题（本大题共8小题，共72分。解答应写出文字说明、演算步骤或证明过程。）}

\vspace{6pt}

\begin{enumerate}[label=\textbf{\arabic*.}, leftmargin=2.5em, resume]

\item （8分）计算：
\begin{enumerate}[label=(\arabic*), nosep]
\item $(-8) + 5 - (-3) + (-7)$
\item $(-1)^{2023} \times (-2) + 16 \div (-4)$
\end{enumerate}
\vspace{6cm}

\item （8分）化简求值：$3(2a^2b - ab^2) - 2(a^2b - 2ab^2)$，其中 $a = 2$，$b = -1$。
\vspace{6cm}

\item （8分）解方程：
\begin{enumerate}[label=(\arabic*), nosep]
\item $3x - 7 = 8 - 2x$
\item $\dfrac{x-1}{3} - \dfrac{x+2}{6} = 1$
\end{enumerate}
\vspace{6cm}

\item （9分）某文具店购进笔记本和钢笔共 $100$ 支（本），笔记本每本进价 $5$ 元，钢笔每支进价 $10$ 元，共花费 $650$ 元。
\begin{enumerate}[label=(\arabic*), nosep]
\item 笔记本和钢笔各购进了多少？
\item 若笔记本售价 $8$ 元，钢笔售价 $15$ 元，全部售出后可获利多少元？
\end{enumerate}
\vspace{6cm}

\item （9分）如图，线段 $AB = 12$，$C$ 是 $AB$ 的中点，$D$ 是 $CB$ 上一点，且 $DB = 4$。
\begin{enumerate}[label=(\arabic*), nosep]
\item 求 $AC$ 的长。
\item 求 $CD$ 的长。
\end{enumerate}
\vspace{6cm}

\item （10分）已知 $|a - 2| + (b + 3)^2 = 0$。
\begin{enumerate}[label=(\arabic*), nosep]
\item 求 $a$、$b$ 的值。
\item 求 $a^2 - 2ab + b^2$ 的值。
\end{enumerate}
\vspace{6cm}

\item （10分）如图，已知 $\angle AOB = 90^{\circ}$，$\angle BOC = 40^{\circ}$，$OM$ 平分 $\angle AOC$。
\begin{enumerate}[label=(\arabic*), nosep]
\item 求 $\angle AOC$ 的度数。
\item 求 $\angle AOM$ 的度数。
\item 求 $\angle BOM$ 的度数。
\end{enumerate}
\vspace{6cm}

\item （10分）将连续的正整数 $1, 2, 3, \cdots$ 按照如下规律排列成数阵：
\begin{center}
\begin{tabular}{ccccc}
$1$ & $2$ & $4$ & $7$ & $\cdots$ \\
$3$ & $5$ & $8$ & $\cdots$ & \\
$6$ & $9$ & $\cdots$ & & \\
$10$ & $\cdots$ & & & \\
$\vdots$ & & & &
\end{tabular}
\end{center}
观察以上排列规律，回答下列问题：
\begin{enumerate}[label=(\arabic*), nosep]
\item 第 $3$ 行第 $1$ 个数是 $6$，第 $4$ 行第 $1$ 个数是 $10$，则第 $5$ 行第 $1$ 个数是多少？
\item 第 $n$ 行第 $1$ 个数是多少（用含 $n$ 的代数式表示）？
\item 数 $28$ 在第几行第几列？
\end{enumerate}
\vspace{6cm}

\end{enumerate}

\end{document}
